\documentclass[11pt,a4paper]{article}
\usepackage[margin=27mm,headheight=15pt]{geometry}
\usepackage[T1]{fontenc}
\usepackage[utf8]{inputenc}
\usepackage{lmodern}
\usepackage{microtype}
\usepackage{amsmath,amssymb,amsthm,mathtools}
\usepackage{booktabs,tabularx,array}
\usepackage{enumitem}
\usepackage{fancyhdr}
\usepackage{xurl}
\usepackage[hidelinks]{hyperref}
\usepackage[nameinlink,noabbrev]{cleveref}
\setlength{\emergencystretch}{2em}
\setlist[enumerate]{itemsep=0.3em,topsep=0.4em}
\setlist[itemize]{itemsep=0.25em,topsep=0.4em}
\setcounter{tocdepth}{2}
\pagestyle{fancy}
\fancyhf{}
\fancyhead[L]{\small Simultaneous convolution divisors}
\fancyhead[R]{\small ProveIt research companion}
\fancyfoot[C]{\thepage}
\newtheorem{theorem}{Theorem}[section]
\newtheorem{proposition}[theorem]{Proposition}
\newtheorem{lemma}[theorem]{Lemma}
\newtheorem{corollary}[theorem]{Corollary}
\newtheorem{definition}[theorem]{Definition}
\newtheorem{example}[theorem]{Example}
\theoremstyle{remark}
\newtheorem{remark}[theorem]{Remark}
\newcommand{\R}{\mathbb R}
\newcommand{\C}{\mathbb C}
\newcommand{\N}{\mathbb N}
\newcommand{\Nzero}{\mathbb N_0}
\newcommand{\Z}{\mathbb Z}
\newcommand{\T}{\mathbb R/\mathbb Z}
\newcommand{\E}{\mathbb E}
\newcommand{\Prob}{\mathbb P}
\newcommand{\Law}{\mathcal L}
\newcommand{\supp}{\operatorname{supp}}
\newcommand{\sinc}{\operatorname{sinc}}
\newcommand{\ord}{\operatorname{ord}}
\newcommand{\lcm}{\operatorname{lcm}}
\newcommand{\Unif}{\mathsf U}
\newcommand{\vp}{v_p}
\newcommand{\ind}{\mathbf 1}
\newcommand{\dil}[1]{\mathsf D_{#1}}
\newcommand{\cf}[1]{\widehat{#1}}
\newcommand{\eqd}{\overset{d}{=}}
\newcommand{\pathtext}[1]{\nolinkurl{#1}}
\DeclareMathOperator{\Var}{Var}
\DeclareMathOperator{\diam}{diam}
\hypersetup{pdftitle={Simultaneous Convolution Divisors of Fabius-Type Laws},
 pdfauthor={Research report prepared for the ProveIt project},
 pdfsubject={Prime-adic factorization, finite certificates, linear rigidity, and residual regularity}}

\title{\bfseries Simultaneous Convolution Divisors\\of Fabius-Type Laws\\[0.5em]
\large A prime-adic Hall criterion, finite certificates,\\linear rigidity, and a sharp regularity transition}
\author{Research report prepared for the ProveIt project}
\date{29 September 2026}

\begin{document}
\maketitle
\begin{abstract}
Let $X_p=\sum_{k\ge0}p^{-k}U_k$, where $p$ is prime and the $U_k$ are independent
uniform random variables on $[-1,1]$. We classify every convolution factor of
$\Law(X_p)$ that is itself a finite or countable convolution of centered uniform
laws. The factor $\Law(\sum_j a_jU_j)$ exists if and only if $a_j=1/n_j$ for positive
integers $n_j$ and
\[
 \#\{j:v_p(n_j)\le K\}\le K+1\qquad(K\ge0).
\]
For this class, cancellation of all poles in the quotient of entire characteristic
functions is already sufficient for positivity. A nested matching argument gives
an explicit independent residual and finite obstruction certificates. For finitely
many geometric streams, the infinitely many inequalities reduce to one finite
window and a rational density condition.

In several dimensions, every uniform generator of a convolution factor of an
independent product law must lie on a coordinate axis. Consequently, its complete
linear decomposability semigroup consists of partial monomial matrices whose
nonzero entries are signed reciprocal integers. The scalar matching theorem
extends to prime-power bases, but not to arbitrary composite bases: an explicit
base-six quotient is entire and is the transform of a signed, not positive, measure.

For pure prime-power scales, the residual has a canonical digit representation
whose multiplicities are exactly the unused matching capacities. Positive capacity
density gives a compactly supported smooth density, a sharp log-squared Fourier
limsup, and differential non-finiteness. At zero capacity density, an eventually
periodic digit profile gives a complete classification into finite atomic laws,
splines, singular continuous laws, or densities of an explicitly determined exact
finite differentiability order. All mathematical assertions designated as theorems
are proved below; a standard-library Python companion checks finite certificates
and examples.
\end{abstract}

\paragraph{Status and scope.}
This is a proof-bearing research manuscript, not a report of a Lean verification or
peer review. It develops extensions of the repository's already existing
reciprocal-integer \emph{single-copy} divisor theorem, rather than claiming that
result as new. The source review was targeted, not an exhaustive audit of every
repository artifact or the entire literature. Global priority for the extensions
proved here is not asserted.

\clearpage
\tableofcontents
\clearpage

\section{The question and its relationship to the repository}
\label{sec:intro}

The Fabius and Rvachev functions provide a particularly concrete meeting point of
probability, entire functions, arithmetic, and formal mathematics. Their Fourier
transforms are infinite products of sinc functions. A zero of such a product is
not just a point where a transform vanishes: its multiplicity records how many
geometric scales coincide at that frequency. The question addressed here is
whether that arithmetic record completely determines which independent uniform
components can be removed from the distribution.

The existing ProveIt report \emph{Reciprocal-Integer Convolution Divisors of the
Rvachev Law} already constructs and classifies decompositions of the form
$X\eqd cX'+Y$, where $X'$ is an independent copy of the same Rvachev law
\cite{RepoDivisors}. Its reciprocal-integer classification, quotient cocycle, and
odd/even regularity results are prior work relative to this manuscript. The
repository also contains exact arithmetic descriptions of the zero divisor in
\pathtext{GeneralizedZeroDivisor.lean} \cite{RepoLean}. Those results suggest a
strictly larger problem:
\begin{quote}
Which \emph{collections} of scaled uniform laws, or of geometric uniform
convolutions with different bases, can occur simultaneously as independent
components of a fixed Fabius-type law?
\end{quote}
The distinction is substantial. Two factors can each be individually admissible
but impossible to remove together. Conversely, several different geometric
streams can interleave exactly, or leave a singular Cantor residual rather than a
smooth function. Scalar variance inequalities do not detect these phenomena.

The main contribution is an exact answer for prime and prime-power target bases.
The answer has four interchangeable forms: convolution divisibility, analytic
zero cancellation, a family of arithmetic inequalities, and a constructive
matching of source factors to target scales. The matching is a special nested
instance of the representative-selection principle associated with Hall's theorem
\cite{Hall}; its proof here requires only sorting integer deadlines.

\subsection{What is established here, and what is reused}

\begin{center}
\begin{tabularx}{\textwidth}{@{}p{0.29\textwidth}X@{}}
\toprule
Part of the argument & Role and status \\
\midrule
Infinite sinc product and its zeros & Established background; proved again in the normalization used here. \\
Single scaled independent copy & Existing repository result; recovered as a corollary, not claimed as a new contribution. \\
Simultaneous uniform factors & Complete classification and explicit residual construction in \cref{sec:hall}. \\
Geometric streams & Necessary integer-base restriction and a complete finite decision criterion in \cref{sec:streams}. \\
Multidimensional factors & Coordinate-direction rigidity and the complete matrix classification in \cref{sec:multi}. \\
Analytic versus positive factorization & Equivalence for prime-power targets and a separating base-six example in \cref{sec:primepower}. \\
Residual regularity & Exact subcritical Fourier limsup and critical periodic classification in \cref{sec:smooth,sec:critical}. \\
\bottomrule
\end{tabularx}
\end{center}

The terminology of a linear decomposability or Urbanik semigroup is conventional;
Jurek describes its relation to generalized Mehler semigroups \cite{Jurek}.
The matrix classification below is proved directly and does not invoke an
external semigroup-classification theorem.

\subsection{A normalization warning}

We use uniforms on $[-1,1]$ and the characteristic-function convention
$\cf\mu(t)=\int e^{itx}\,d\mu(x)$. Thus $X_2$ is supported on $[-2,2]$.
The classical Rvachev up-density on $[-1,1]$ is the density of $X_2/2$.
Under the repository's convention $\int e^{2\pi izx}\,d\mu(x)$, the transform of
$X_2/2$ is $\prod_{k\ge0}\sinc(\pi z/2^k)$, as in the classical account
\cite{Arias}. Every scalar decomposition theorem is unchanged by this common
normalization of source and target; explicit supports and digit locations are not.

\section{Probability and analytic preliminaries}
\label{sec:prelim}

For $a>0$, let $\Unif_a$ denote the uniform law on $[-a,a]$, and put
\[
 \sinc z=\begin{cases}\sin z/z,&z\ne0,\\1,&z=0.\end{cases}
\]
If $b>1$ is real and the $U_k$ are independent with law $\Unif_1$, define
\begin{equation}
 X_b=\sum_{k=0}^{\infty}b^{-k}U_k,
 \qquad \mu_b=\Law(X_b),
 \qquad \Phi_b(z)=\prod_{k=0}^{\infty}\sinc(z/b^k).
 \label{eq:base-law}
\end{equation}
The series converges absolutely for every outcome. Its support is the interval
$[-b/(b-1),b/(b-1)]$: finite sums have interval support with radii tending to
$b/(b-1)$, and the remaining summands have radii tending to zero. Independence
and bounded convergence give $\cf{\mu_b}(t)=\Phi_b(t)$. In particular,
\begin{equation}
 \E X_b=0,
 \qquad \Var(X_b)=\frac13\sum_{k\ge0}b^{-2k}
                 =\frac{b^2}{3(b^2-1)}.
 \label{eq:variance}
\end{equation}

For a finite or countable indexed family $a=(a_j)$ of positive numbers with
$\sum_j a_j<\infty$, let
\begin{equation}
 \nu_a=\Law\left(\sum_j a_jU_j\right),
 \qquad \Psi_a(z)=\prod_j\sinc(a_jz).
 \label{eq:uniform-source}
\end{equation}
The empty family gives the point mass $\delta_0$ and the product $1$.
All copies of uniforms or other random variables occurring in an asserted
independent sum are understood to be mutually independent.

\begin{definition}[Convolution divisibility]
A probability law $\nu$ divides a probability law $\mu$, written $\nu\mid_*\mu$,
if there is a probability law $\rho$ such that $\mu=\nu*\rho$.
The residual $\rho$ is not assumed to belong to any prescribed family.
For a scalar $c$, $\dil c\nu$ denotes the pushforward of $\nu$ by $x\mapsto cx$.
\end{definition}

\begin{lemma}[Entire products and exact zeros]
\label{lem:zeros}
The products in \eqref{eq:base-law} and \eqref{eq:uniform-source} converge locally
uniformly on $\C$ and are entire. They have precisely the zeros contributed by
the individual factors, with multiplicities added. For prime $p$,
\begin{equation}
 Z(\Phi_p)=\{\pi m:m\in\Z\setminus\{0\}\},
 \qquad \ord_{\pi m}\Phi_p=1+v_p(|m|).
 \label{eq:target-zeros}
\end{equation}
Here $v_p(n)$ is the exponent of $p$ in a positive integer $n$.
\end{lemma}
\begin{proof}
On a fixed compact set, $\sinc(a_jz)=1+O(a_j^2)$ uniformly once $a_j$ is small.
Since $\sum a_j^2<\infty$, the product converges locally uniformly. Away from
zeros of the finitely many large factors, a sufficiently far tail admits a
uniformly absolutely convergent logarithm and hence is nonzero. At any fixed
nonzero $z$, only finitely many factors can vanish, because a vanishing factor
requires $a_j|z|\ge\pi$. The zeros of every individual sinc factor are simple.
These observations prove the first assertion and the multiplicity rule.

For $a_k=p^{-k}$, a zero is $\pi p^k\ell$ with $\ell$ a nonzero integer.
Their union is exactly the nonzero integer multiples of $\pi$. At $\pi m$,
the vanishing factors are those with $p^k\mid m$, namely
$k=0,\ldots,v_p(|m|)$.
\end{proof}

\begin{lemma}[Compact residuals and uniqueness]
\label{lem:residual}
If compactly supported probability laws satisfy $\mu=\nu*\rho$, then $\rho$ is
compactly supported. If $\cf\nu$ extends to a nonzero entire function, the
residual is unique. In particular, $\nu_a\mid_*\mu_b$ implies that
$\Phi_b/\Psi_a$ has an entire extension.
\end{lemma}
\begin{proof}
For $x\in\supp\nu$ and $y\in\supp\rho$, every sufficiently small neighborhood
of $x+y$ receives positive convolution mass from neighborhoods of $x$ and $y$.
Thus $x+y\in\supp\mu$. Fixing one $x\in\supp\nu$ confines $\supp\rho$ to a
translate of the compact set $\supp\mu$. The characteristic function of a
compactly supported law is entire, by differentiation under a bounded-support
integral. The convolution identity extends from the real line to $\C$ by the
identity theorem, proving the quotient assertion.

The real zeros of a nonzero entire function are isolated. Hence two residual
characteristic functions agree on a dense subset of $\R$ by cancellation, and
then everywhere by continuity. Uniqueness of characteristic functions gives
uniqueness of the probability laws.
\end{proof}

\section{A complete prime-adic Hall criterion}
\label{sec:hall}

Each target uniform $p^{-k}U_k$ will be called supply slot $k$. A source uniform
$U/n$ can be extracted from that slot whenever $p^k\mid n$. The crucial feature
is that the eligible slots are a nested initial interval, not an arbitrary set.

\begin{lemma}[Splitting one uniform]
\label{lem:split}
If $n,r\in\N$ and $a=r/n$, then
\begin{equation}
 \Unif_a=\Unif_{1/n}*\Law(V_{n,r}),\qquad
 \Prob\left(V_{n,r}=\frac{2\ell-r+1}{n}\right)=\frac1r
 \quad(0\le\ell<r).
 \label{eq:split}
\end{equation}
For $r=1$ the digit law is $\delta_0$.
\end{lemma}
\begin{proof}
The $r$ translates of $[-1/n,1/n]$ centered at the displayed points tile
$[-r/n,r/n]$, with overlaps only at endpoints. Each mixture component has
weight $1/r$ and density $n/2$, so the combined density is $1/(2a)$ on that
interval and zero outside it. Notice also that $|V_{n,r}|\le a-1/n$.
\end{proof}

\begin{theorem}[Complete simultaneous-factor classification]
\label{thm:hall}
Let $p$ be prime and let $\nu_a$ be as in \eqref{eq:uniform-source}. The following
conditions are equivalent.
\begin{enumerate}[label=\textup{(\roman*)}]
\item $\nu_a\mid_*\mu_p$.
\item $\Phi_p/\Psi_a$ extends to an entire function on $\C$.
\item There are positive integers $n_j$ such that $a_j=1/n_j$ and
\begin{equation}
 D(K):=\#\{j:v_p(n_j)\le K\}\le K+1
 \qquad\text{for every }K\in\Nzero.
 \label{eq:hall}
\end{equation}
\item There are positive integers $n_j=1/a_j$ and an injection
$j\mapsto k_j$ into $\Nzero$ such that $p^{k_j}\mid n_j$ for every $j$.
\end{enumerate}
When these conditions hold, the residual is unique and has an explicit
representation as an absolutely convergent sum of independent finite digit
variables and the unmatched target uniforms.
\end{theorem}
\begin{proof}
\emph{(i)$\Rightarrow$(ii).} This is \cref{lem:residual}.

\emph{(ii)$\Rightarrow$(iii).} The first positive zero $\pi/a_j$ of any source
factor must be a zero of $\Phi_p$. By \eqref{eq:target-zeros}, $1/a_j$ is a
positive integer, denoted $n_j$.

Suppose \eqref{eq:hall} fails at $K$. Choose $K+2$ source indices whose
$v_p(n_j)$ are at most $K$ and put $m$ equal to the least common multiple of
these denominators. Then $v_p(m)\le K$. Every chosen factor vanishes at $\pi m$,
so $\Psi_a$ has order at least $K+2$ there, whereas $\Phi_p$ has order at most
$K+1$. The quotient has a pole, a contradiction. This argument also handles
an infinite set of demands at a fixed deadline: a finite subfamily suffices.

\emph{(iii)$\Rightarrow$(iv).} Write $d_j=v_p(n_j)$. Condition
\eqref{eq:hall} makes each bounded-deadline set finite. Therefore the source
indices can be enumerated in nondecreasing deadline order, with arbitrary
fixed ordering of ties. Let $d_{(k)}$ be the deadline of the zero-indexed
$k$th source in this enumeration. If $d_{(k)}<k$, the first $k+1$ sources
would all have deadline at most $d_{(k)}$, contradicting
$D(d_{(k)})\le d_{(k)}+1$. Thus $d_{(k)}\ge k$.
Assign the $k$th source to slot $k$. This is the required injection.
For a finite family only the corresponding finite initial set of slots is used.

\emph{(iv)$\Rightarrow$(iii).} Every source of deadline at most $K$ must be
assigned to one of slots $0,\ldots,K$. Injectivity gives \eqref{eq:hall}.

\emph{(iv)$\Rightarrow$(i).} For the source assigned to slot $k_j$, put
$r_j=n_j/p^{k_j}\in\N$. By \cref{lem:split},
\[
 p^{-k_j}U\eqd n_j^{-1}U'+V_{n_j,r_j}.
\]
Use this splitting independently at every matched slot. Keep the entire uniform
$p^{-k}U_k$ at every unmatched slot. All the residual summands have absolute
value at most $p^{-k}$ at their respective slots, so their sum converges
absolutely. The extracted source variables also converge absolutely. Their
independence and the slot-by-slot identity give
$\mu_p=\nu_a*\rho$. Uniqueness follows from \cref{lem:residual}.
\end{proof}

\begin{remark}[Why positivity is a conclusion]
An entire quotient of two characteristic functions need not be a characteristic
function. In \cref{thm:hall}, positivity follows from the explicit independent
construction, not from the analytic cancellation by itself. The arithmetic
structure allows the cancellation condition to force a valid construction.
Section~\ref{sec:primepower} gives an entire quotient for which positivity fails
when that structure is absent.
\end{remark}

\begin{remark}[Summability is also forced by the certificate]
If a countable family of positive integers satisfies \eqref{eq:hall}, its sorted
enumeration has $n_{(k)}\ge p^k$. Consequently
$\sum_j1/n_j\le p/(p-1)$. Thus the arithmetic condition itself supplies the
absolute-convergence assumption needed to define the source probability law.
\end{remark}

\subsection{Finite witnesses and exact residual bounds}

\begin{corollary}[A finite certificate for every failure]
\label{cor:witness}
For reciprocal-integer widths, failure of convolution divisibility has a finite
witness: a finite subfamily whose least common multiple $m$ satisfies
\[
 \#\{j\text{ in the subfamily}:n_j\mid m\}>1+v_p(m).
\]
For a finite input family, sorting deadlines and testing $d_{(k)}\ge k$ either
constructs the residual slot matching or produces such a witness.
\end{corollary}
\begin{proof}
If $d_{(k)}<k$, take the first $k+1$ sorted sources. The valuation of their
least common multiple is $d_{(k)}$, so the source order exceeds the target
order. Conversely, such an excess prevents an entire quotient by
\cref{lem:zeros}. The algorithm terminates after the finite sort and scan.
\end{proof}

\begin{proposition}[Support endpoints, variance, and a truncation certificate]
\label{prop:residual-bounds}
For an admissible source, put
\[
 R=\frac{p}{p-1}-\sum_j\frac1{n_j}.
\]
The residual is symmetric, is supported in $[-R,R]$, and has both $-R$ and $R$
in its support. Its variance is
\begin{equation}
 \Var(\rho)=\frac13\left(\frac{p^2}{p^2-1}-\sum_j\frac1{n_j^2}\right).
 \label{eq:residual-variance}
\end{equation}
If $\rho^{[K]}$ retains only residual summands assigned to target slots
$0,\ldots,K-1$, then
\begin{equation}
 W_\infty(\rho,\rho^{[K]})\le\frac{p^{1-K}}{p-1}.
 \label{eq:winfty}
\end{equation}
Here $W_\infty$ is the infimum of the essential maximum displacement over all
couplings. No assertion that the residual support is an interval is intended.
\end{proposition}
\begin{proof}
Each split digit law is symmetric and has extreme support points
$\pm(p^{-k_j}-n_j^{-1})$. Unmatched slots have extreme points $\pm p^{-k}$.
Summing these radii gives $R$. Every finite set of summands can simultaneously
lie arbitrarily near their upper endpoints with positive probability, and the
remaining total radius tends to zero. Thus $R$, and similarly $-R$, belongs to
the support. Variances subtract in the independent decomposition and yield
\eqref{eq:residual-variance}. Finally, couple the full residual and its truncation
using the same retained summands. The discarded absolute displacement is at
most $\sum_{k\ge K}p^{-k}=p^{1-K}/(p-1)$.
\end{proof}

\section{Geometric streams: an infinite problem with a finite certificate}
\label{sec:streams}

We next allow finitely many entire Fabius-type laws as source components. A
geometric stream means the sequence of uniforms making up one scaled copy of
$X_b$. Besides these streams, finitely many individual uniform factors may be
present.

\begin{theorem}[Classification of geometric source families]
\label{thm:streams}
Fix prime $p$. Consider the law of
\begin{equation}
 Y=\sum_{\ell=1}^{J}u_\ell U_\ell+
   \sum_{i=1}^{r}c_iX_{b_i}^{(i)},
 \qquad u_\ell,c_i>0,\quad b_i>1,
 \label{eq:streams}
\end{equation}
with all components independent. It divides $\mu_p$ if and only if all of the
following hold:
\begin{enumerate}[label=\textup{(\alph*)}]
\item $u_\ell=1/N_\ell$ and $c_i=1/n_i$ for positive integers $N_\ell,n_i$;
\item each $b_i$ is an integer divisible by $p$;
\item writing $d_\ell=v_p(N_\ell)$, $a_i=v_p(n_i)$, and
$s_i=v_p(b_i)\ge1$, one has
\begin{equation}
 D(K)=\sum_{\ell=1}^{J}\ind_{d_\ell\le K}
      +\sum_{i=1}^{r}\max\left\{0,1+\left\lfloor\frac{K-a_i}{s_i}\right\rfloor\right\}
      \le K+1
 \quad(K\ge0).
 \label{eq:stream-demand}
\end{equation}
\end{enumerate}
Negative nonzero coefficients are handled by replacing them with their absolute
values, since all source laws are symmetric. Zero coefficients may be discarded.
\end{theorem}
\begin{proof}
Expand every $X_{b_i}$ into its uniform summands and apply \cref{thm:hall}.
The first term of each stream forces $c_i=1/n_i$, and all stream denominators
$n_ib_i^j$ must be integers. The case $j=1$ makes $b_i$ rational. Write
$b_i=u/v$ in lowest terms. Since $n_iu^j/v^j$ is an integer for every $j$,
$v^j$ divides $n_i$ for every $j$, and therefore $v=1$. Thus $b_i$ is an integer.
If $p$ did not divide $b_i$, the infinitely many denominators of this stream
would all have the same $p$-adic valuation, violating \eqref{eq:hall}.

For an integer base divisible by $p$, the deadlines of stream $i$ are exactly
$a_i+s_ij$ for $j\ge0$. Counting how many are at most $K$ gives
\eqref{eq:stream-demand}. The converse follows by the same expansion and the
sufficiency part of \cref{thm:hall}.
\end{proof}

Define
\begin{equation}
 A=\max(\{0\}\cup\{d_\ell\}\cup\{a_i\}),\quad
 L=\lcm(s_1,\ldots,s_r),\quad
 \eta=\sum_{i=1}^{r}\frac1{s_i},\quad \delta=1-\eta.
 \label{eq:density}
\end{equation}
For no streams, the conventions are $L=1$ and $\eta=0$. Put
$F(K)=K+1-D(K)$. The parameter $\delta$ is the asymptotic surplus of supply
slots over source demands per valuation level.

\begin{theorem}[Finite-window criterion]
\label{thm:finite-window}
All inequalities in \eqref{eq:stream-demand} hold if and only if
\begin{equation}
 \delta\ge0\quad\text{and}\quad
 F(K)\ge0\quad\text{for }0\le K<A+L.
 \label{eq:finite-window}
\end{equation}
More precisely, for every $K\ge A$,
\begin{equation}
 F(K+L)=F(K)+L\delta,
 \qquad L\delta=L-\sum_iL/s_i\in\Z.
 \label{eq:drift}
\end{equation}
\end{theorem}
\begin{proof}
For $K\ge A$, every finite demand has already arrived and every stream has
started. Increasing $K$ by $L$ adds exactly $L/s_i$ demands from stream $i$ and
adds $L$ supply slots. This proves \eqref{eq:drift}. If $L\delta<0$, iterating
along any residue class eventually makes $F$ negative. If $L\delta\ge0$,
every $K\ge A$ is obtained from a unique element of $[A,A+L-1]$ by adding a
nonnegative multiple of $L$, and its slack cannot decrease. The earlier
indices are precisely the remaining finite checks in \eqref{eq:finite-window}.
\end{proof}

\paragraph{Algorithmic interpretation.}
Compute the valuations, $A$, $L$, and the integer drift $L\delta$. Scan the
$A+L$ displayed inequalities. A negative entry is a finite rejection
certificate. If all pass but the drift is negative, let
$q=\lfloor F(A)/(-L\delta)\rfloor+1$; then $K=A+qL$ is an explicit violating
index. If all pass and the drift is nonnegative, \eqref{eq:drift} certifies
all remaining indices. The method is exact, but its least-common-multiple
window need not be polynomial in the binary length of the input. An implementation
must distinguish a resource limit from a negative mathematical certificate.

\begin{example}[Density is necessary but not sufficient]
For $p=2$, two unscaled copies of $X_4$ have deadline streams $0,2,4,\ldots$.
Their combined density is $1/2+1/2=1$, but they already demand two slots at
$K=0$, where only one is available. In contrast, the two streams
$0,2,4,\ldots$ and $1,3,5,\ldots$ use every slot once and give
\[
 X_2\eqd X_4+\tfrac12X_4'.
\]
The simultaneous criterion distinguishes these cases before any analytic
approximation or numerical positivity test is needed.
\end{example}

\begin{corollary}[The existing single-copy theorem, recovered]
\label{cor:single}
For $c\ne0$, $\dil c\mu_p\mid_*\mu_p$ if and only if $|c|=1/n$ for a positive
integer $n$. Including $c=0$ gives the trivial factor $\delta_0$.
\end{corollary}
\begin{proof}
A single stream with base $p$ has step $s=1$ and start $a=v_p(n)\ge0$;
its demand is $\max(0,K-a+1)\le K+1$. Necessity is already part of
\cref{thm:streams}. For $p=2$, this recovers the classification established
in the repository's divisor report \cite{RepoDivisors}.
\end{proof}

\section{Multidimensional rigidity and the matrix semigroup}
\label{sec:multi}

The scalar classification has a strong geometric consequence. Let $d\ge1$ and
consider the target product law $\mu_p^{\otimes d}$ on $\R^d$. Its characteristic
function is
\[
 \Phi_{p,d}(t)=\prod_{q=1}^{d}\Phi_p(t_q).
\]
Its real zero set is a countable union of coordinate hyperplanes. A source
uniform along an oblique line creates an oblique zero hyperplane, and there is
no positive residual that can remove it.

\begin{theorem}[Rigidity of uniform generators]
\label{thm:vector}
Let $v_j\in\R^d\setminus\{0\}$ be a finite or countable family with
$\sum_j\|v_j\|<\infty$, and let $Y=\sum_jv_jU_j$. Then
$\Law(Y)\mid_*\mu_p^{\otimes d}$ if and only if every generator has the form
\begin{equation}
 v_j=\varepsilon_j\frac{e_{q(j)}}{n_j},\qquad
 \varepsilon_j\in\{-1,1\},\quad n_j\in\N,
 \label{eq:vector-form}
\end{equation}
and each coordinate separately satisfies
\begin{equation}
 \#\{j:q(j)=q,\ v_p(n_j)\le K\}\le K+1
 \qquad(q=1,\ldots,d;\ K\ge0).
 \label{eq:row-hall}
\end{equation}
The residual is unique and is a product of its coordinate residual laws.
\end{theorem}
\begin{proof}
Suppose first that the factorization exists. Fix a generator $v=v_j$. The
source characteristic function vanishes everywhere on
$H=\{t\in\R^d:v\cdot t=\pi\}$. Thus the target must vanish there, and
\[
 H\subseteq\bigcup_{q=1}^{d}\bigcup_{m\in\Z\setminus\{0\}}
                  \{t:t_q=\pi m\}.
\]
If $v$ has at least two nonzero coordinates, no coordinate hyperplane in the
union contains $H$. Its intersection with $H$ is consequently either empty
or a proper affine subspace of $H$, hence closed and nowhere dense in $H$.
The affine space $H$ is complete and cannot be a countable union of such
subsets by the Baire category theorem. This contradiction proves that every
$v_j$ lies on a coordinate axis. When $d=1$, that conclusion is automatic.

After grouping generators by coordinate and absorbing their signs by symmetry,
the coordinate projections of the purported factorization are scalar
factorizations of $\mu_p$. Apply \cref{thm:hall} to each projection. It gives
\eqref{eq:vector-form} and \eqref{eq:row-hall}.

Conversely, the scalar construction supplies a residual in every coordinate.
Taking their independent product gives the required vector decomposition.
Finally, the source characteristic function is nonzero on a dense subset of
$\R^d$: it is a product of one-variable entire functions with isolated real
zeros. Thus cancellation and continuity make the residual characteristic
function unique. It must be the product already constructed.
\end{proof}

\begin{definition}[Linear decomposability semigroup]
For a probability law $\mu$ on $\R^d$, define
\[
 \mathcal D(\mu)=\{C\in\R^{d\times d}:\mu=(C_*\mu)*\rho
                    \text{ for some probability law }\rho\}.
\]
This is a semigroup under matrix multiplication: combining factorizations for
$C$ and $D$ gives a factorization for $CD$. A partial monomial matrix has at most
one nonzero entry in each row and at most one in each column.
\end{definition}

\begin{theorem}[Complete matrix classification]
\label{thm:matrix}
For prime $p$ and every $d\ge1$,
\begin{equation}
 \mathcal D(\mu_p^{\otimes d})=
 \left\{C:\ C\text{ is partial monomial and every nonzero entry is }
                    \frac{\varepsilon}{n},\ 
                    \varepsilon\in\{-1,1\},\ n\in\N\right\}.
 \label{eq:matrix-semigroup}
\end{equation}
The invertible matrices in this set are signed permutation matrices with
reciprocal-integer coordinate scalings. Its semigroup units are exactly the
signed permutation matrices. Every continuous path taking values in
$\mathcal D(\mu_p^{\otimes d})$ is constant.
\end{theorem}
\begin{proof}
Let $X'$ have independent coordinates with law $\mu_p$. Writing each coordinate
as its uniform series expresses $CX'$ using the generators $p^{-k}c_j$, where
$c_j$ is column $j$ of $C$. By \cref{thm:vector}, every nonzero column is
coordinate-aligned and its coefficient is a signed reciprocal integer.

If two nonzero columns were assigned to the same row, that row would contain
two deadline streams, each of step $1$. Their combined demand is $2K+O(1)$,
which eventually exceeds $K+1$. Hence each row contains at most one nonzero
column, proving the necessary partial monomial structure. Conversely, every
nonzero row of such a matrix invokes exactly one admissible scalar
single-copy factor from \cref{cor:single}; zero rows leave a full target
coordinate as residual. The independent rowwise constructions prove sufficiency.

Invertibility forces exactly one nonzero entry in each row and column. The
inverse can also satisfy \eqref{eq:matrix-semigroup} only when every scaling
integer is $1$, proving the assertion about units. Finally, each matrix entry
of a continuous path lies in the countable subset
$\{0\}\cup\{\pm1/n:n\in\N\}$ of $\R$. A connected subset of $\R$ with more
than one point contains an interval and is uncountable. Each entry, and hence
the path, must therefore be constant.
\end{proof}

\begin{corollary}[No continuous nontrivial linear scaling family]
There is no continuous nonconstant family $C_t\in\mathcal D(\mu_p^{\otimes d})$
with $C_0=I$. In particular, no continuous linear contraction semigroup tending
to zero can provide decompositions of this fixed target law at every time.
\end{corollary}
\begin{proof}
Apply the final assertion of \cref{thm:matrix}.
\end{proof}

This is a rigidity statement about \emph{exact independent decompositions}; it
is not a prohibition on approximate decompositions, nonlinear dynamics, or
unrelated stochastic processes. It also explains why the scalar reciprocal
integers should not be mistaken for a continuous self-decomposability property.
With different source bases, several columns may feed the same output
coordinate, but the corresponding row must then satisfy the stream criterion
of \cref{thm:streams}.

\section{Prime-power targets and a composite-base obstruction}
\label{sec:primepower}

\begin{theorem}[Prime-power extension]
\label{thm:primepower}
Let $B=p^h$ with $p$ prime and $h\ge1$. For a summable uniform source $\nu_a$,
convolution divisibility $\nu_a\mid_*\mu_B$ and entire extensibility of
$\Phi_B/\Psi_a$ are equivalent. They hold exactly when $a_j=1/n_j$ with
$n_j\in\N$ and
\begin{equation}
 \#\left\{j:\left\lfloor\frac{v_p(n_j)}h\right\rfloor\le K\right\}
 \le K+1\qquad(K\ge0).
 \label{eq:primepower-hall}
\end{equation}
The residual construction assigns source $j$ to a distinct slot $k_j$ with
$B^{k_j}\mid n_j$ and splits that uniform as in \cref{lem:split}.
\end{theorem}
\begin{proof}
The target zeros are still precisely $\pi m$, $m\ne0$, but their orders are
\[
 1+\left\lfloor\frac{v_p(|m|)}h\right\rfloor.
\]
The allowed slots of denominator $n$ are the initial interval through
$\lfloor v_p(n)/h\rfloor$. The least-common-multiple obstruction works because
\[
 \left\lfloor\frac{v_p(\lcm(n_1,\ldots,n_s))}h\right\rfloor
 =\max_i\left\lfloor\frac{v_p(n_i)}h\right\rfloor.
\]
The necessity, sorted matching, and independent splitting proofs of
\cref{thm:hall} therefore apply word for word with these deadlines and
supply scales $B^{-k}$.
\end{proof}

For a general integer $B$, the eligible supply slots of a single denominator
are still nested. What fails is the least-common-multiple reduction: different
prime factors can be supplied by different denominators and can create extra
powers of $B$ in their least common multiple. The resulting analytic zero
condition can miss a genuine positivity obstruction.

\begin{theorem}[An entire nonpositive quotient at base six]
\label{thm:base6}
The function
\begin{equation}
 Q(t)=\frac{\Phi_6(t)}{\sinc(t/2)\sinc(t/3)}
 \label{eq:base6-quotient}
\end{equation}
extends to an entire function, is normalized by $Q(0)=1$, and is not the
characteristic function of a probability measure. In particular,
$\Unif_{1/2}*\Unif_{1/3}$ does not divide $\mu_6$.
\end{theorem}
\begin{proof}
Peeling off the first two factors of $\Phi_6$ and using elementary trigonometric
identities gives
\begin{align}
 Q(t)
 &=\frac{\sinc t\,\sinc(t/6)}{\sinc(t/2)\sinc(t/3)}\Phi_6(t/36)\notag\\
 &=\frac{\cos(t/2)}{\cos(t/6)}\Phi_6(t/36)\notag\\
 &=\bigl(2\cos(t/3)-1\bigr)\Phi_6(t/36).
 \label{eq:base6-signed}
\end{align}
The identities first hold away from their displayed denominators' zeros and
then extend analytically. The final expression is entire.

Let $\tau=\dil{1/36}\mu_6$. Its support is
$[-1/30,1/30]$. Equation~\eqref{eq:base6-signed} is the characteristic transform
of the finite signed measure
\[
 (\delta_{-1/3}+\delta_{1/3}-\delta_0)*\tau.
\]
The three translates have disjoint containing intervals. On the central
interval $[-1/30,1/30]$, this signed measure has mass exactly $-1$.
It is therefore not positive. Uniqueness of Fourier transforms for finite
signed measures prevents a different positive measure from having the same
transform. Finally, an actual convolution residual in
\eqref{eq:base6-quotient} would have characteristic function $Q$, contradicting
what was just proved.
\end{proof}

The example has an analytic quotient with every pole canceled, so it cannot be
rejected by searching for an unremoved Fourier zero. It shows precisely why a
positivity construction is indispensable and why the prime-power theorem is
not a routine assertion about arbitrary entire quotients.

\section{The residual digit profile}
\label{sec:digits}

From now until explicitly stated otherwise, the target base is prime $p$ and
the source denominators are pure powers $n_j=p^{d_j}$ with $d_j\in\Nzero$.
Assume the Hall inequalities hold. Define the \emph{digit slack}
\begin{equation}
 g_h=h-\#\{j:d_j<h\}=F(h-1),\qquad h\ge1.
 \label{eq:digit-slack}
\end{equation}
In particular, $0\le g_h\le h$. This is the number of target digit layers
remaining at level $h$ after all source digit layers are removed.

Let $D_p$ be uniform on the centered digit set
$\{2\ell-p+1:0\le\ell<p\}$ and write
\begin{equation}
 \chi_p(u)=\E e^{iuD_p}
 =\frac1p\sum_{\ell=0}^{p-1}e^{i(2\ell-p+1)u}
 =\frac{\sin(pu)}{p\sin u}.
 \label{eq:digit-char}
\end{equation}
The quotient has its removable values at multiples of $\pi$; the finite sum
is its globally valid definition.

\begin{lemma}[Uniform digit expansion]
\label{lem:uniform-digits}
For every complex $z$,
\begin{equation}
 \sinc z=\prod_{h=1}^{\infty}\chi_p(z/p^h).
 \label{eq:uniform-digits}
\end{equation}
Equivalently, $\sum_{h\ge1}p^{-h}D_{p,h}$ has law $\Unif_1$ for independent
copies $D_{p,h}$ of $D_p$.
\end{lemma}
\begin{proof}
The finite product telescopes to
$\sin z/[p^N\sin(z/p^N)]$, with removable values understood. Its locally
uniform limit is $\sinc z$. The digit series is absolutely convergent, and
its characteristic function is the product on the right. Uniqueness of
characteristic functions proves the probability statement.
\end{proof}

\begin{theorem}[Canonical residual digit law]
\label{thm:digits}
For the pure-power source above, the unique residual has the representation
\begin{equation}
 Z=\sum_{h=1}^{\infty}p^{-h}\sum_{\ell=1}^{g_h}D_{h,\ell},
 \qquad
 \cf\rho(t)=\prod_{h=1}^{\infty}\chi_p(t/p^h)^{g_h},
 \label{eq:residual-digits}
\end{equation}
where the $D_{h,\ell}$ are independent with law $D_p$.
The series is absolutely convergent. Moreover,
\begin{equation}
 \ord_{\pi p^K}\cf\rho=g_{K+1}=F(K)\qquad(K\ge0).
 \label{eq:slack-zeros}
\end{equation}
\end{theorem}
\begin{proof}
Expand every uniform using \cref{lem:uniform-digits}. The target product has
exponent $h$ at digit level $h$: the pairs consisting of supply scale $k\ge0$
and a positive relative digit index with sum $h$ number exactly $h$.
A source uniform $p^{-d_j}U$ contributes one digit at every level $h>d_j$.
The difference of the two multiplicities is exactly \eqref{eq:digit-slack}.

Since $\sum_h h p^{-h}<\infty$ and $|D_p|\le p-1$, all these series and
regroupings are absolutely convergent. The source digits plus the digits
in \eqref{eq:residual-digits} reproduce the target digits independently.
This proves the factorization and hence, by uniqueness, the claimed residual.
The zero-order formula also follows by subtracting the source order
$\#\{j:d_j\le K\}$ from the target order $K+1$ at $\pi p^K$.
\end{proof}

The construction in \cref{thm:hall} is a matching of \emph{whole uniforms};
\cref{thm:digits} instead matches their \emph{individual digit layers}. They
produce the same law but answer different questions. The first is well suited
to an exact existence certificate for arbitrary integer denominators. The
second exposes regularity, arithmetic scaling, and entire-function structure.

For a finite-plus-geometric pure-power source, the definitions in
\eqref{eq:density} give
\begin{equation}
 g_h=\delta h+O(1).
 \label{eq:slack-asymptotic}
\end{equation}
If $\delta=0$, \eqref{eq:drift} says that $g_h$ is periodic with period $L$
for every $h\ge A+1$. These two regimes lead to very different residuals.

\section{Positive surplus: smoothness, sharp Fourier decay, and non-D-finiteness}
\label{sec:smooth}

\begin{theorem}[A sharp log-squared Fourier envelope]
\label{thm:sharp}
Let a pure-power source satisfy the Hall condition and suppose its digit slack
has $g_h=\delta h+O(1)$ with $\delta>0$. Then its residual has a compactly
supported $C^\infty$ density and
\begin{equation}
 \log|\cf\rho(t)|\le
 -\frac{\delta}{2\log p}(\log|t|)^2+O(\log|t|)
 \qquad(|t|\to\infty),
 \label{eq:fourier-upper}
\end{equation}
where $\log0=-\infty$. The coefficient is sharp:
\begin{equation}
 \limsup_{|t|\to\infty}
 \frac{\log|\cf\rho(t)|}{(\log|t|)^2}
 =-\frac{\delta}{2\log p}.
 \label{eq:fourier-limsup}
\end{equation}
The limsup is attained as a limit along $t_N=\pi p^N/(p+1)$.
\end{theorem}
\begin{proof}
Choose a nonnegative constant $C$ so that $g_h\ge\delta h-C$ for all $h$,
and define the nondecreasing integer minorant
\[
 e_h=\max(0,\lfloor\delta h-C\rfloor)\le g_h.
\]
For $j\ge1$, let $H_j$ be the first index with $e_h\ge j$. Then
$H_j=j/\delta+O(1)$, and another application of the uniform digit expansion
gives
\begin{equation}
 \prod_{h\ge1}\chi_p(t/p^h)^{e_h}
   =\prod_{j\ge1}\sinc\bigl(t/p^{H_j-1}\bigr).
 \label{eq:minorant-uniforms}
\end{equation}
The remaining digit multiplicities $g_h-e_h$ define a probability law, so their
characteristic function has modulus at most $1$ on the real axis.

Put $T=\log_p|t|$ and retain from \eqref{eq:minorant-uniforms} the
$J=\delta T+O(1)$ factors with $H_j-1\le T$. The elementary bound
$|\sinc u|\le\min(1,1/|u|)$ yields
\begin{align*}
 \log|\cf\rho(t)|
 &\le -J\log|t|+\log p\sum_{j=1}^{J}(H_j-1)\\
 &=-J\log|t|+\frac{\log p}{2\delta}J^2+O(J)\\
 &=-\frac{\delta}{2\log p}(\log|t|)^2+O(\log|t|).
\end{align*}
This proves the upper bound uniformly in real $t$ of large modulus, including
zeros with the stated convention.

For sharpness, put $\alpha=\pi/(p+1)$ and $t_N=\alpha p^N$. Since
$p\equiv-1\pmod{p+1}$, for every integer $j\ge0$,
\[
 |\sin(\alpha p^j)|=\sin\alpha,
 \qquad |\chi_p(\alpha p^j)|=1/p.
\]
The first $N$ digit factors in \eqref{eq:residual-digits} therefore contribute
\[
 -\log p\sum_{h=1}^{N}g_h
 =-\tfrac12\delta\log p\,N^2+O(N)
\]
to the logarithmic modulus. The remaining arguments are
$\alpha/p^r$, $r\ge1$, on an interval where $\chi_p$ is positive and nonzero.
Taylor's theorem and compactness give
$0\le-\log\chi_p(\alpha/p^r)\le C_1p^{-2r}$ for a fixed constant $C_1$.
Since $g_{N+r}=O(N+r)$, the entire tail contributes $O(N)$ in absolute value.
Thus
\[
 \log|\cf\rho(t_N)|=-\tfrac12\delta\log p\,N^2+O(N).
\]
Dividing by $(N\log p+\log\alpha)^2$ proves the lower bound in
\eqref{eq:fourier-limsup}; the upper bound was already established.

Finally, \eqref{eq:fourier-upper} implies
$\int_\R |t|^m|\cf\rho(t)|\,dt<\infty$ for every nonnegative integer $m$.
Fourier inversion supplies a continuous density, and differentiation under the
integral supplies continuous derivatives of every order. The law was already
compactly supported, so this density is compactly supported as well.
\end{proof}

\begin{corollary}[Smoothness for arbitrary integer stream denominators]
\label{cor:general-smooth}
Every admissible finite-plus-geometric family in \cref{thm:streams} with
$\delta>0$ has a compactly supported $C^\infty$ residual density. Its
characteristic function satisfies the upper bound \eqref{eq:fourier-upper}
with the same $\delta$. Sharpness of that coefficient is asserted here only
for the pure-power case of \cref{thm:sharp}.
\end{corollary}
\begin{proof}
Replace each source denominator $n$ by $p^{v_p(n)}$, making that source uniform
larger. The new family has exactly the same deadlines, is admissible by
\cref{thm:hall}, and has a pure-power residual $\rho_0$ governed by
\cref{thm:sharp}. Each enlarged source uniform splits into the original
uniform and an independent finite digit law, because $n/p^{v_p(n)}$ is an
integer. Summing these splits is legitimate: the enlarged widths are summable
by the Hall bound. Thus the enlarged source equals the original source
convolved with a probability law $\sigma$, and the original residual is
$\rho_0*\sigma$. Its Fourier modulus is bounded by $|\cf{\rho_0}|$.
Fourier inversion gives the stated smoothness and upper bound.
\end{proof}

\begin{definition}[D-finiteness]
An entire function $f$ is D-finite if it satisfies a nontrivial homogeneous
linear differential equation
$\sum_{j=0}^{r}P_j(z)f^{(j)}(z)=0$ with polynomial coefficients $P_j$.
\end{definition}

\begin{theorem}[Positive-surplus residuals are not D-finite]
\label{thm:nonDfinite}
For every admissible finite-plus-geometric source from \cref{thm:streams}
with $\delta>0$, the entire residual characteristic function is not D-finite.
This includes non-pure integer bases and denominators.
\end{theorem}
\begin{proof}
At valuation cutoff $K$, the source denominators with $p$-adic valuation at
most $K$ form a finite set of size $D(K)$. Let $q_K$ be the least common
multiple of their prime-to-$p$ parts, with the empty least common multiple
equal to $1$, and put $M_K=p^Kq_K$. Exactly those $D(K)$ source factors vanish
at $\pi M_K$: their denominators divide $M_K$, while every other denominator
has $p$-adic valuation greater than $K$. The target order there is $K+1$.
Consequently
\[
 \ord_{\pi M_K}\cf\rho=K+1-D(K)=\delta K+O(1),
\]
which tends to infinity. These are infinitely many distinct points since
$v_p(M_K)=K$.

Suppose a differential equation of order $r$ with leading coefficient
$P_r\not\equiv0$ existed. A nonzero entire function cannot satisfy a nontrivial
order-zero equation, so $r\ge1$. At any point where $P_r$ does not vanish,
local uniqueness for the linear analytic differential equation implies that a
solution with its first $r$ initial derivatives zero is identically zero.
Thus a nonzero solution has zero multiplicity less than $r$ at every such
ordinary point. The leading polynomial has only finitely many exceptional
points, contradicting the infinite sequence of unbounded multiplicities.
The residual is nonzero because $\cf\rho(0)=1$.
\end{proof}

The zero-multiplicity obstruction for the original Rvachev transform is already
present in the repository's representation work \cite{RepoAudit}. The content
here is its extension to the residuals, where the remaining zero order is the
explicit matching slack and is controlled by the geometric demand density.

\section{Critical periodic slack: a complete regularity classification}
\label{sec:critical}

Assume the pure-power source is admissible and its slack $g_h$ is eventually
periodic. This includes every finite-plus-geometric pure-power family with
$\delta=0$, by \cref{thm:finite-window}. The result below is more general than
that application: it only needs eventual periodicity of the actual slack.

The distinction between an atomless singular law and a discrete law matters.
Convolving either with a uniform produces an absolutely continuous law, but the
singular continuous case gains one more order of continuous differentiability
than the discrete case. We first prove the precise singularity statement needed
for this distinction.

\begin{lemma}[A missing periodic digit layer forces singular continuity]
\label{lem:periodic-singular}
Let $(r_h)_{h\ge1}$ be a periodic sequence of nonnegative integers. Suppose
$\min_h r_h=0$ and the sequence is not identically zero. Then
\begin{equation}
 S=\sum_{h\ge1}p^{-h}\sum_{\ell=1}^{r_h}D_{h,\ell}
 \label{eq:periodic-singular}
\end{equation}
has a singular continuous law: it is singular with respect to Lebesgue measure
and has no atoms.
\end{lemma}
\begin{proof}
Replace each centered digit $D_{h,\ell}$ by $2K_{h,\ell}-(p-1)$, where
$K_{h,\ell}$ is uniform on $\{0,\ldots,p-1\}$. This changes the law by an
invertible affine transformation. It therefore suffices to study
\[
 W=\sum_{h\ge1}p^{-h}\sum_{\ell=1}^{r_h}K_{h,\ell}.
\]
Let $L$ be a period and $B=p^L$. Grouping consecutive blocks of $L$ levels gives
\begin{equation}
 W=\sum_{j=0}^{\infty}B^{-(j+1)}A_j,
 \label{eq:block-digits}
\end{equation}
where the $A_j$ are independent, identically distributed, integer-valued finite
digits. They need not lie between $0$ and $B-1$; that restriction is not used.
Modulo one, multiplication by $B$ deletes the initial block. Consequently the
projected law $\lambda$ on $\T$ is invariant under $T_B(x)=Bx\pmod1$ and is a
factor of the independent block shift.

The independent block shift is ergodic. Indeed, sufficiently separated cylinder
events are independent; approximation by cylinder events gives mixing for all
measurable events. For a shift-invariant event $E$, this gives
$\Prob(E)=\Prob(E)^2$. Every invariant event under the factor map pulls back
to such an event, so $\lambda$ is ergodic as well.

Choose $r_0\in\{1,\ldots,L\}$ with $r_{r_0}=0$, and set $n=p^{r_0-1}$.
The Fourier coefficient $\int e^{2\pi inx}\,d\lambda(x)$ is a product of digit
coefficients. At levels $h<r_0$ all factors are $1$. The potentially vanishing
factor at level $h=r_0$ is absent. At levels $h>r_0$, the digit ratio is
$e^{2\pi i/p^{h-r_0+1}}$, so its average over $p$ consecutive digits is
nonzero. The resulting infinite tail product is nonzero as well: the digit
multiplicities are bounded and the differences of these factors from $1$ are
summable. Thus $\lambda$ has a nonzero Fourier coefficient at the nonzero
integer $n$, and is not Haar measure.

For completeness, an ergodic $T_B$-invariant probability on $\T$ that is not
Haar must be singular. The map $T_B$ sends absolutely continuous measures to
absolutely continuous measures and singular measures to singular measures:
for the latter assertion, a null support remains null under the piecewise
linear map. Uniqueness of the Lebesgue decomposition therefore makes the
absolutely continuous and singular components of an invariant law separately
invariant. Ergodicity forces one component to have full mass. One way to see
this last assertion directly is to let $f=d\lambda_{\rm ac}/d\lambda$, so
$0\le f\le1$. Invariance of both measures gives
\[
 \int f(f\circ T_B)\,d\lambda=\int f^2\,d\lambda,
 \qquad
 \int(f\circ T_B)^2\,d\lambda=\int f^2\,d\lambda.
\]
Hence $f=f\circ T_B$ almost everywhere, and ergodicity makes $f$ constant.
If that constant is positive, the whole measure is absolutely continuous.
But an absolutely continuous invariant probability has Fourier coefficients
$\cf\lambda(n)=\cf\lambda(Bn)$, which tend to zero for $n\ne0$ by the
Riemann--Lebesgue lemma. It is therefore Haar. Since the present law is not
Haar, it is singular. The inverse image in $\R$ of a null set of full
$\lambda$-measure is a countable union of null translates, proving singularity
of the real law in \eqref{eq:block-digits}.

Finally, there are infinitely many levels with $r_h>0$. Select one raw digit
at each such level and put zeros at all other positions. The resulting base-$p$
random expansion is atomless: every real number has at most two base-$p$
expansions with digits $0,\ldots,p-1$, and every fixed infinite choice at the
selected positions has probability zero. The remaining digits are independent
of this selected subseries. Convolution with an atomless measure remains
atomless, since the mass of a singleton is the integral of singleton masses.
This proves the assertion for $W$ and hence for $S$.
\end{proof}

\begin{lemma}[A compact probability has no vanishing positive-order derivative]
\label{lem:nonzero-derivative}
If $\rho$ is a compactly supported probability measure and $m\ge1$, its
$m$th distributional derivative is not the zero distribution.
\end{lemma}
\begin{proof}
If it vanished, its Fourier transform would give
$(-it)^m\cf\rho(t)=0$ for all real $t$ in our sign convention. Thus
$\cf\rho(t)=0$ for $t\ne0$, contradicting continuity and $\cf\rho(0)=1$.
\end{proof}

\begin{theorem}[Exact critical regularity]
\label{thm:critical}
Let an admissible pure-power source have slack $g_h$ periodic for $h\ge H$.
Set
\[
 m=\min_{h\ge H}g_h,\qquad r_h=g_h-m\quad(h\ge H),\qquad a=p^{1-H}.
\]
Let $P$ be the finite atomic law of the prefix digit sum at levels $h<H$.
Then the residual has the factorization
\begin{equation}
 \rho=P*\Unif_a^{*m}*\sigma,
 \label{eq:critical-decomposition}
\end{equation}
where $\sigma$ is the tail law made from the $r_h$ digits. Empty convolution
powers mean $\delta_0$. Exactly the following alternatives hold.
\begin{enumerate}[label=\textup{(\roman*)}]
\item If $r_h=0$ throughout the periodic tail and $m=0$, then $\rho=P$ is
finite atomic.
\item If $r_h=0$ throughout the tail and $m\ge1$, then $\rho$ is a finite
mixture of compact splines. For $m=1$ it has a piecewise constant density but
no continuous density. For $m\ge2$ its density belongs to $C^{m-2}(\R)$ and
not to $C^{m-1}(\R)$.
\item If the periodic tail $r_h$ is not identically zero and $m=0$, then
$\rho$ is singular continuous.
\item If the periodic tail $r_h$ is not identically zero and $m\ge1$, then
$\rho$ has a density in $C^{m-1}(\R)$ and not in $C^m(\R)$.
\end{enumerate}
These regularity statements concern existence of a globally continuous
representative of the density and its derivatives, not arbitrary changes of
a density on null sets.
\end{theorem}
\begin{proof}
Separate the finite prefix in \cref{thm:digits}, and remove $m$ copies of a
complete infinite digit tail. By \cref{lem:uniform-digits}, each complete tail
starting at $H$ is uniform on $[-p^{1-H},p^{1-H}]$. This proves
\eqref{eq:critical-decomposition}. If $r$ is nonzero, its periodic minimum is
zero, and rescaling and reindexing the tail puts it under
\cref{lem:periodic-singular}. Thus $\sigma$ is singular continuous.

The distributional derivative of the uniform measure is
\begin{equation}
 D\Unif_a=\frac{\delta_{-a}-\delta_a}{2a}=:B_a.
 \label{eq:uniform-derivative}
\end{equation}
If $\sigma=\delta_0$, the residual is a finite sum of translates of an
$m$-fold uniform convolution. For $m\ge2$, each distributional derivative
through order $m-2$ is a finite sum of translates of a convolution retaining
at least two uniforms, and hence has a continuous density. This gives a
$C^{m-2}$ density. Its derivative of order $m-1$ is a compactly supported
piecewise constant function, because
$D^{m-1}\rho=P*B_a^{*(m-1)}*\Unif_a$.
If the original density were $C^{m-1}$, this piecewise constant derivative
would have a continuous representative and would therefore be identically
zero. That contradicts \cref{lem:nonzero-derivative}. The same compact
piecewise-constant argument applied to the density itself handles $m=1$.
The case $m=0$ is immediate.

Now suppose $\sigma$ is singular continuous. Convolution with the finite
atomic $P$ preserves both singularity and absence of atoms, proving (iii).
For $m\ge1$, the law $\Unif_a*\sigma$ has the continuous density
\[
 x\longmapsto\frac{F_\sigma(x+a)-F_\sigma(x-a)}{2a},
\]
where $F_\sigma$ is continuous because $\sigma$ is atomless. Equation
\eqref{eq:uniform-derivative} shows that the distributional derivative of
$\rho$ through order $m-1$ has a continuous density, so $\rho$ has a
$C^{m-1}$ density. Its next derivative is the finite signed measure
\[
 D^m\rho=P*B_a^{*m}*\sigma.
\]
This measure is singular, being a finite signed sum of translates of a
singular measure, and is nonzero by \cref{lem:nonzero-derivative}. A $C^m$
density would make the same derivative absolutely continuous, which is
impossible. This proves (iv).
\end{proof}

\paragraph{Why the classification is effective.}
For finitely many pure-power streams, the certificate already supplies a
period $L$ and a valid start $H=A+1$. Computing the $L$ tail values of $g_h$
determines their minimum $m$ and whether all values equal that minimum.
The four-way classification thus needs only integer arithmetic after the
existence test. A shorter actual period may be used but is not required.
In contrast, no exact critical-regularity classification for arbitrary
prime-to-$p$ parts is claimed here.

\section{Explicit factorizations and diagnostic examples}
\label{sec:examples}

\subsection{Two identical source streams and a Cantor residual}

\begin{theorem}[A singular residual from two identical geometric factors]
\label{thm:cantor-example}
For every prime $p$,
\begin{equation}
 X_p\eqd\frac1pX_{p^2}^{(1)}+\frac1pX_{p^2}^{(2)}+Z_p,
 \qquad
 Z_p=\sum_{j=0}^{\infty}p^{-(2j+1)}D_j,
 \label{eq:cantor-decomposition}
\end{equation}
with independent terms and independent centered $p$-digit variables $D_j$.
The residual $Z_p$ is singular continuous and has Hausdorff dimension $1/2$.
Its extreme support points and variance are
\begin{equation}
 \pm\frac{p}{p+1},\qquad
 \Var(Z_p)=\frac{p^2}{3(p^2+1)}.
 \label{eq:cantor-stats}
\end{equation}
\end{theorem}
\begin{proof}
Both source streams have start $1$ and step $2$. Their slack is
$g_h=1$ at odd levels and $0$ at even levels, which satisfies the Hall
inequalities. The canonical digit formula is exactly
\eqref{eq:cantor-decomposition}; singular continuity follows from
\cref{lem:periodic-singular}. The extreme radius is
\[
 (p-1)\sum_{j\ge0}p^{-(2j+1)}=\frac{p}{p+1}.
\]
Since $\Var(D_j)=(p^2-1)/3$, summing the independent variances gives
\eqref{eq:cantor-stats}.

The residual obeys the iterated-function identity
$Z_p\eqd D_0/p+Z_p'/p^2$. Its $p$ first-level support intervals are disjoint
with positive gaps, and every cylinder of depth $N$ has diameter
$[2p/(p+1)]p^{-2N}$ and mass $p^{-N}$. Covering by these cylinders shows that
the support has Hausdorff dimension at most $1/2$. Conversely, for any small
radius $r$, choose $N$ with $p^{-2N}$ comparable to $r$. An interval of
length $r$ meets at most a fixed number of depth-$N$ cylinders, by their
uniformly separated geometry. Its mass is at most a constant times
$p^{-N}\le C r^{1/2}$. Summing this bound over arbitrary interval covers
shows the dimension is at least $1/2$. The same bound and the cylinder
masses give dimension $1/2$ for the probability measure itself.
\end{proof}

For $p=2$, $D_j$ is a fair sign, and the formula becomes
\[
 X_2\eqd\tfrac12X_4^{(1)}+\tfrac12X_4^{(2)}
             +\sum_{j\ge0}\frac{\varepsilon_j}{2\cdot4^j}.
\]
The residual is a centered, rescaled middle-half Cantor measure with support
endpoints $\pm2/3$ and variance $4/15$. Each source component has a smooth
density, but after both are removed the remaining law is singular. Thus the
regularity of an admissible residual is not inferred from the regularity of
its source components.

\subsection{Every finite differentiability order from two base-four copies}

\begin{corollary}[A tunable exact regularity family]
\label{cor:tunable}
For every integer $a\ge1$, two independent copies of $2^{-a}X_4$ jointly divide
$X_2$. If $a=1$, their residual is the Cantor law above. If $a\ge2$, the residual
has a density in $C^{a-2}(\R)$ and not in $C^{a-1}(\R)$.
\end{corollary}
\begin{proof}
The two deadline streams both have start $a$ and step $2$. Their slack begins
$g_h=h$ for $1\le h\le a$, then alternates between $a-1$ and $a$.
It is nonnegative, so the factorization exists. The periodic tail has minimum
$m=a-1$ and is not constant. Apply \cref{thm:critical}.
\end{proof}

In particular, two quarter-scaled base-four laws leave a continuous but not
continuously differentiable density. Here the exact decomposition is
\[
 \rho=\Unif_1*\Law\left(\sum_{j\ge1}4^{-j}\varepsilon_j\right).
\]
Two eighth-scaled copies leave a $C^1$ but not $C^2$ density. These exact endpoint
regularity obstructions arise from nonzero singular distributional derivatives,
not merely from an inconclusive upper bound on a Fourier transform.

\subsection{Checks across all four critical regimes}

\begin{center}
\begin{tabularx}{\textwidth}{@{}p{0.32\textwidth}p{0.22\textwidth}X@{}}
\toprule
Source removed from $X_2$ & First digit slacks & Residual conclusion \\
\midrule
$X_4+X_4'/2$ & $0,0,0,0,\ldots$ & Point mass at zero. \\
$X_2'/2+U/2$ & $1,0,0,0,\ldots$ & Fair atoms at $\pm1/2$. \\
$X_2'/2$ & $1,1,1,1,\ldots$ & Uniform on $[-1,1]$; no continuous density. \\
$X_2'/4$ & $1,2,2,2,\ldots$ & A continuous piecewise linear density, not $C^1$. \\
$(X_4^{(1)}+X_4^{(2)})/2$ & $1,0,1,0,\ldots$ & Singular continuous Cantor law. \\
$(X_4^{(1)}+X_4^{(2)})/4$ & $1,2,1,2,\ldots$ & $C^0$ density, not $C^1$. \\
$(X_4^{(1)}+X_4^{(2)})/8$ & $1,2,3,2,3,\ldots$ & $C^1$ density, not $C^2$. \\
\bottomrule
\end{tabularx}
\end{center}

The finite-spline rows recover familiar finite-head decompositions, and are
included as normalization and boundary-case checks rather than as novelty
claims. The singular and mixed smooth--singular rows demonstrate the additional
behavior introduced by simultaneous streams.

\subsection{A subcritical check and a rejected critical case}

Removing one unscaled $X_4$ from $X_2$ leaves $\tfrac12X_4'$. Its stream has
$\eta=1/2$, so $\delta=1/2$; the sharp coefficient in
\eqref{eq:fourier-limsup} is $-1/(4\log2)$, agreeing with the direct base-four
product. On the other hand, two unscaled copies of $X_4$ fail at $K=0$ despite
$\eta=1$. Their characteristic-function quotient has a pole at $\pi$: the
source order there is $2$ and the target order is $1$. This is the smallest
possible exact rejection certificate.

\subsection{Why periodicity cannot be dropped from the critical theorem}

\begin{example}[Zero surplus density with a smooth residual]
\label{ex:squares}
Remove from $X_p$ exactly the target uniforms whose supply exponents are not
perfect squares. The residual is then
\[
 Z=\sum_{j=0}^{\infty}p^{-j^2}U_j.
\]
The removed source has asymptotic deadline density $1$, because the number of
unremoved exponents through $K$ is only $\lfloor\sqrt K\rfloor+1$.
Nevertheless, the residual is $C^\infty$. For any fixed $N$, its first $N$
uniform factors give a Fourier bound $O(|t|^{-N})$, while all remaining
characteristic factors have modulus at most $1$. Taking $N$ arbitrarily large
and applying Fourier inversion proves smoothness. Its slack is unbounded and
nonperiodic. Thus the exact classification in \cref{thm:critical} depends on
periodicity, not on the bare statement that the asymptotic surplus density is
zero.
\end{example}

\section{Exact computation and a formalization boundary}
\label{sec:verification}

\subsection{The accompanying executable certificates}

The file \pathtext{verify.py} uses Python 3.10 or later and only the standard
library. Running
\begin{quote}\ttfamily
python verify.py --output verification.json
\end{quote}
reproduces the exact checks and the example certificates. Its public routines
have the following roles.

\begin{center}
\begin{tabularx}{\textwidth}{@{}p{0.31\textwidth}X@{}}
\toprule
Routine & Contract \\
\midrule
\pathtext{finite_certificate} & For target base $p^{\mathtt{power}}$, returns a sorted slot matching or a least-common-multiple zero-order obstruction. \\
\pathtext{stream_certificate} & Returns a complete finite-window certificate, or a violating deadline including its demand and capacity. Inputs are valuation starts, steps, and optional finite deadlines. \\
\pathtext{demand} & Computes the exact integer expression in \eqref{eq:stream-demand}. \\
\pathtext{run_tests} & Compares certificates with independent finite matching and zero-order checks, verifies moment identities, and writes diagnostic examples. \\
\bottomrule
\end{tabularx}
\end{center}

The distinction between a source denominator and its valuation is explicit:
\pathtext{finite_certificate} takes denominators, whereas
\pathtext{stream_certificate} takes deadlines. Invalid primes, nonpositive
denominators, invalid steps, and exceeded resource limits raise exceptions.
In particular, a large least common multiple is not silently treated as a
mathematical rejection.

The execution accompanying this article produced the following results.
\begin{center}
\begin{tabular}{@{}lr@{}}
\toprule
Exact check & Number completed \\
\midrule
Finite uniform families & 9,100 \\
Finite families accepted & 546 \\
Uniform interval partitions & 96 \\
Rational moment identities through degree ten & 1,056 \\
Finite-plus-geometric families & 8,855 \\
Geometric families accepted & 1,275 \\
Integer period-drift identities & 153,390 \\
Base-six Laurent-polynomial identity & 1 \\
Deliberately invalid or resource-limited cases & 4 \\
\bottomrule
\end{tabular}
\end{center}

The finite-family tests cover all multisets of length at most four with
denominators in $\{1,\ldots,12\}$ for target bases $2,3,5,4,9$.
The stream tests cover all multisets of at most three streams with starts
$0,\ldots,4$ and steps $1,\ldots,4$, each combined with five finite-deadline
patterns. The finite matching oracle uses backtracking rather than the sorted
deadline routine. The zero-order oracle checks least common multiples of
subfamilies. The stream oracle checks a longer direct window and verifies every
returned rejection witness, while separate exact tests check the period recurrence.
All uniform moment calculations use rational arithmetic.

These computations are finite implementation tests, not a replacement for the
infinite arguments in the article. They do not numerically establish positivity,
regularity, or a limit theorem. Those conclusions come from the proofs.

\subsection{Existing Lean infrastructure and what remains to be formalized}

The inspected file \pathtext{GeneralizedZeroDivisor.lean} already supplies a
precise dyadic zero-divisor interface for weighted products. In particular,
the source contains the declarations
\begin{quote}\small\ttfamily
Fabius.generalizedRvachevProduct\_eq\_zero\_iff\_int\\
Fabius.analyticOrderAt\_generalizedRvachevProduct\_int
\end{quote}
with the summability and nonzero-integer hypotheses stated in that file
\cite{RepoLean}. The first identifies the arithmetic zero set; the second
identifies its multiplicity. The report-local README also explicitly distinguishes
these adjacent zero-divisor theorems from a formalization of the divisor
classification itself \cite{RepoDivisorsReadme}.

This manuscript did not compile the repository or invoke Lean. It therefore
makes no claim that its new theorems are already kernel checked. A realistic
formalization sequence would be:
\begin{enumerate}
\item Establish the finite nested matching equivalence between prefix counts
and sorted deadlines, including a proof-producing rejection witness.
\item Extend that matching to countably many demands by the finite cardinality
of every bounded-deadline set; record the bound on the sum of matched widths.
\item Formalize the one-uniform digit split and its independent countable
product construction, so that positivity is obtained from probability semantics.
\item Connect analytic zero order to the least-common-multiple obstruction,
using the existing dyadic interface first and then general prime powers.
\item Prove the geometric deadline recurrence and verify the finite certificate
checker against its universal quantified specification.
\item Formalize the vector zero-hyperplane obstruction and the rowwise reduction.
The regularity and ergodic arguments should be a separate final development,
not hidden as assumptions inside the factorization theorem.
\end{enumerate}
The main existence theorem can be formalized without the critical singularity
or smoothness theory. Keeping that dependency boundary explicit makes an initial
verified implementation smaller and its trust story clearer.

\section{Further research questions and proposed projects}
\label{sec:research}

The following are open directions \emph{within the development of this article}.
They are not assertions that no corresponding theorem exists elsewhere in the
literature. Each question specifies what the present results do and do not settle.

\begin{enumerate}[label=\textbf{Q\arabic*.},leftmargin=3em]
\item \textbf{Positive factorization for non-prime-power targets.}
For integer $B$ with at least two distinct prime divisors, characterize all
uniform-convolution factors of $\mu_B$. The base-six example proves that zero
multiplicity is insufficient. A promising exact formulation should combine
prime-vector divisibility with positivity of finite cyclotomic quotient masks,
rather than simply replacing $v_p$ by the largest exponent of $B$ dividing an
integer. The first useful milestone is a complete decision procedure for a fixed
small number of source uniforms.

\item \textbf{The full convolution-factor monoid.}
The main theorem classifies factors that are themselves sums of independent
uniforms. It does not classify every compact probability factor of $\mu_p$.
The digit representation exhibits many purely discrete and singular factors
outside that input class. Is there a canonical factor description in terms of
integer-valued digit multiplicities and finite positive masks, and what
indecomposable factors must be added? An answer would distinguish a genuine
factorization theorem for all laws from the uniform-source theorem proved here.

\item \textbf{Critical streams with prime-to-$p$ parts.}
Existence depends only on valuations, but the residual law remembers the full
integers. Determine exact smoothness, singularity, and Fourier dimension when
$\delta=0$ and some source denominators or bases have nontrivial prime-to-$p$
parts. The pure-power residual construction remains a convolution factor, but
that alone does not determine the effect of the additional digit measures.
The repository's one-stream reciprocal-integer family is an important test case,
not a substitute for the simultaneous classification.

\item \textbf{Nonperiodic capacity profiles.}
Find necessary and sufficient regularity conditions directly in terms of the
admissible sequence $g_h$. Example~\ref{ex:squares} shows that zero asymptotic
surplus permits a smooth law when the leftover scales are sparse and nonperiodic.
Can a bounded but nonperiodic profile produce a smooth residual? Which growth
conditions give prescribed sub-log-squared Fourier scales, and what replaces the
explicit maximizing sequence in \cref{thm:sharp}?

\item \textbf{Shorter arithmetic certificates.}
The finite-window algorithm is exact but can have a large period
$\lcm(s_1,\ldots,s_r)$. Seek a certificate whose size is polynomial, or at least
subexponential, in the binary input length for natural restricted stream classes.
Useful possibilities include residue-class compression, rational generating
functions for demand, and short proofs that a floor-sum never exceeds its affine
capacity. Any speedup should still return a witness when the condition fails.

\item \textbf{Stability under approximate factorization.}
Exact reciprocal integrality and coordinate alignment are brittle conclusions.
Quantify the distance from a proposed source to the admissible set when only an
approximate convolution identity is known in a specified metric. A useful result
must state its dependence on source support, observation bandwidth, and the
smallest relevant Fourier zero. In several dimensions, a particularly concrete
goal is a lower bound on approximation error caused by a small but nonzero
oblique component.

\item \textbf{Other target hyperplane arrangements.}
Replace the independent coordinate target by an infinite box spline whose
uniform directions form a nontrivial finite arrangement. The proof of
\cref{thm:vector} suggests that every source direction must coincide with a target
zero-hyperplane direction, but the multiplicity accounting can couple several
arithmetic chains. Determine when a direction-by-direction matching is sufficient
and when additional positivity obstructions, analogous to base six, appear.

\item \textbf{Dimensions and fine regularity of critical digit laws.}
For periodic excess profiles taking only the values $0$ and $1$, separated digit
constructions give accessible dimensions. Repeated digit layers introduce carries
and overlaps. Compute exact dimensions and local dimensions for these profiles,
and determine fractional H\"older, Sobolev, or Besov regularity after the uniform
smoothing layers are added. The present theorem settles exact integer
$C^k$ thresholds but not these finer scales.

\item \textbf{Certified construction and formal verification.}
Implement a proof-producing checker for finite and geometric demands, then
connect its output to a Lean probability-law construction. The natural initial
endpoint is the prime-base existence equivalence, not every analytic corollary.
A later certified sampler could use the uniform residual-tail bound
\eqref{eq:winfty} to return a sample coupled to the exact residual with a rigorous
error radius. This would connect arithmetic certificates, probability semantics,
and practical computation in a single development.
\end{enumerate}

\section{Conclusion}

For prime-base Fabius laws, a simultaneous factorization question in probability
reduces exactly to a nested arithmetic allocation problem. Every source uniform
must have reciprocal-integer width, and the number of sources whose eligible
slots end by level $K$ cannot exceed $K+1$. The least-common-multiple zero
obstruction proves necessity, while a uniform digit split proves sufficiency
and supplies the residual constructively.

The consequences go beyond an existence test. Geometric source families admit
finite certificates. Multidimensional sources cannot use oblique directions,
and exact linear copy-factors form a countable matrix semigroup with no
nonconstant continuous paths.
A prime-power extension and a base-six signed quotient identify the boundary
between an analytic divisor and a positive convolution factor. Finally, unused
digit capacity determines the residual's regularity: positive linear growth
gives sharp log-squared Fourier decay, while periodic critical capacity gives an
exact atomic--spline--singular--smooth classification.

The mathematical core is self-contained and does not assume the correctness of
unverified repository research claims. Its most direct next step is a verified
matching-and-splitting implementation, followed by the harder classification of
positive factors beyond prime-power target bases.

\clearpage
\appendix
\section{Source provenance and the scope of the review}
\label{app:sources}

The repository's main tree was recorded at commit
\begin{quote}\small\ttfamily
23dd71d2d3d5e85d2b97a8cc45f55e735d69a4ae.
\end{quote}
The root and Fabius READMEs, the representation source audit, and the specific
Lean zero-divisor source were inspected at that pinned snapshot. GitHub's
indexed code-search results supplied the existing reciprocal-integer divisor
report and its README at a second explicit snapshot,
\begin{quote}\small\ttfamily
e9c09b12549ea4e3bea7cebc52b450b761f6f581.
\end{quote}
These are two identified source versions, not a claim that all inspected
material came from one synchronized checkout. The mathematical baseline for
the single-copy result is the report cited in \cite{RepoDivisors}.

The review deliberately checked for a collision with the proposed scalar
factorization idea. The existing report was found, and the scalar theorem was
therefore retained only as background. The article's focus shifted to the full
simultaneous uniform-family criterion, multiple geometric streams, the matrix
classification, and the residual capacity profile. The source review does not
establish an exhaustive absence result for these extensions across every current
or historical ProveIt document. Likewise, the external references establish
background and terminology, not a global priority certificate.

The external sources used substantively were Arias de Reyna's account of the
compactly supported smooth function, Hall's original representative-selection
paper, and Jurek's discussion of Urbanik and Mehler semigroups. The proofs here
supply the particular matching, positivity, geometric, and regularity statements
rather than assigning them to those references.

\begin{thebibliography}{9}

\bibitem{Arias}
J. Arias de Reyna,
\emph{An infinitely differentiable function with compact support: Definition and
Properties}, arXiv:1702.05442 (2017), English version of the author's earlier
paper.
\url{https://arxiv.org/abs/1702.05442}.

\bibitem{Hall}
P. Hall,
\emph{On Representatives of Subsets},
Journal of the London Mathematical Society \textbf{s1-10} (1935), no.~1,
26--30. DOI: \href{https://doi.org/10.1112/jlms/s1-10.37.26}{10.1112/jlms/s1-10.37.26}.

\bibitem{Jurek}
Z. J. Jurek,
\emph{On Relations Between Urbanik and Mehler Semigroups},
Probability and Mathematical Statistics \textbf{29} (2009), 297--308.
\url{https://arxiv.org/abs/0811.2989}.

\bibitem{RepoDivisors}
ProveIt research corpus,
\emph{Reciprocal-Integer Convolution Divisors of the Rvachev Law:
Arithmetic semigroups, singular-to-smooth parity, Thue--Morse quotients, zero
zeta functions, and inverse-Fabius approximants}, report dated 30 August 2026.
\href{https://github.com/VladimirReshetnikov/ProveIt/blob/e9c09b12549ea4e3bea7cebc52b450b761f6f581/Analysis/FabiusFunction/docs/semi-formalized-research-frontiers/drafts/spectra-and-arithmetic/Fabius_Rvachev_Reciprocal_Integer_Convolution_Divisors/Fabius_Rvachev_Reciprocal_Integer_Convolution_Divisors.tex}{Pinned repository source (e9c09b12549e)}.

\bibitem{RepoDivisorsReadme}
ProveIt research corpus,
\emph{Reciprocal-Integer Convolution Divisors: README and formalization boundary},
same report directory and snapshot as \cite{RepoDivisors}.
\href{https://github.com/VladimirReshetnikov/ProveIt/blob/e9c09b12549ea4e3bea7cebc52b450b761f6f581/Analysis/FabiusFunction/docs/semi-formalized-research-frontiers/drafts/spectra-and-arithmetic/Fabius_Rvachev_Reciprocal_Integer_Convolution_Divisors/README.txt}{Pinned repository source (e9c09b12549e)}.

\bibitem{RepoLean}
ProveIt contributors,
\emph{GeneralizedZeroDivisor.lean}, in the FabiusFunction Lean development.
\href{https://github.com/VladimirReshetnikov/ProveIt/blob/23dd71d2d3d5e85d2b97a8cc45f55e735d69a4ae/Analysis/FabiusFunction/Lean/FabiusFunction/GeneralizedZeroDivisor.lean}{Pinned repository source (23dd71d2d3d5)}.

\bibitem{RepoAudit}
ProveIt research corpus,
\emph{Source Audit and Novelty-Collision Record}, representation-frontier package.
\href{https://github.com/VladimirReshetnikov/ProveIt/blob/23dd71d2d3d5e85d2b97a8cc45f55e735d69a4ae/Analysis/FabiusFunction/docs/semi-formalized-research-frontiers/drafts/representations/Representation_Frontiers/assets/fabius_rvachev_representation_frontier/SOURCE_AUDIT.md}{Pinned repository source (23dd71d2d3d5)}.

\end{thebibliography}
\end{document}
